\section{Дневник выполнения работы}

На первом этапе была собрана базовая версия программы и выполнен её запуск под утилитой \texttt{gprof}. По отчёту профилирования выяснилось, что значительная доля времени уходит на операции со строками и на обход дерева: в отчёте базовой версии заметны частые вызовы \texttt{std::string::operator==}, \texttt{std::string::operator<} и рекурсивного \texttt{RBTree::\_find}, причём функция \texttt{RBTree::\_find} занимала около половины времени профиля. 

На втором этапе были заменены сравнения строк вида \texttt{==} и \texttt{<} на вызов \texttt{std::string::compare}. После этого число сравнений сократилось, а в профиле исчезла основная нагрузка на комбинацию операторов сравнения, при этом суммарное время выполнения заметно уменьшилось. В оптимизированной версии в профиле уже явно виден прямой вызов \texttt{compare}, а общее время flat profile снизилось с \texttt{0.32} до \texttt{0.17} секунды. 

На третьем этапе методы поиска и вставки были переписаны в итеративном виде. В текущей версии \texttt{RBTree::insert} выполняет обход дерева циклом, используя \texttt{key.compare(current->key)}, а \texttt{RBTree::\_find} также реализован через \texttt{while}. Это позволило убрать накладные расходы рекурсии и дополнительно снизить время работы основных операций. 

После завершения оптимизации программа была проверена на утечки памяти с помощью \texttt{Valgrind}. Отчёт показал, что при завершении работы в памяти не осталось ни одного неосвобождённого блока, а ошибок использования динамической памяти обнаружено не было.

\pagebreak

\section{Выводы о найденных недочётах}

В ходе исследования были выявлены два основных недостатка исходной версии программы.

Во-первых, слишком дорогими оказались множественные сравнения строк. В исходной реализации для одного логического шага часто выполнялось несколько строковых операций, что хорошо видно по профилю: в базовой версии заметны высокие затраты на \texttt{operator==}, \texttt{operator<} и связанные с ними внутренние функции \texttt{std::string}. 

Во-вторых, рекурсивная реализация поиска и вставки добавляла лишние накладные расходы. Профилирование показало, что значительная часть времени уходила не только на балансировку, но и на сам обход дерева. Перевод этих методов в итеративный вид позволил уменьшить число вызовов функций и лучше использовать время процессора. 

При проверке памяти ошибок и утечек выявлено не было, поэтому исправления были сосредоточены именно на производительности, а не на корректности управления памятью.

\pagebreak